\documentclass[11pt]{article}
\usepackage[utf8]{inputenc}
\usepackage[T1]{fontenc}
\usepackage{times}
\usepackage{amsmath,amssymb}
\usepackage{booktabs}
\usepackage{multirow}
\usepackage{graphicx}
\usepackage{xcolor}
\usepackage{hyperref}
\usepackage{url}
\usepackage[margin=1in]{geometry}
\usepackage{natbib}
\usepackage{subcaption}

\newcommand{\gar}{\textsc{GAR}}
\newcommand{\garours}{\textsc{GAR}\textsubscript{ours}}
\newcommand{\gardongre}{\textsc{GAR}\textsubscript{Dongre}}
\newcommand{\tg}{\textsc{ThreadGate}}
\newcommand{\ug}{\textsc{UniformGate}}
\newcommand{\lo}{\textsc{LoRA-only}}
\renewcommand{\ss}{\textsc{SS}}

\title{Why Auxiliary Attention Losses Fail Under Exposure Bias:\\A Multi-Model Diagnostic Study}

\author{Anonymous}
\date{}

\begin{document}
\maketitle

% ============================================================
% ABSTRACT
% ============================================================
\begin{abstract}
Multi-turn dialogue models progressively lose attention to user goals as conversations grow longer. Prior work shows goal information persists in residual streams (probe AUC$=$0.98) even as attention access declines, motivating learnable residual gating. We implement this as \tg{}, a per-layer gating module trained with an auxiliary Goal Accessibility Ratio (\gar{}) loss. \tg{} optimizes \gar{} under teacher-forcing (TF \gar{}$=$0.50) but collapses on model-generated text (AR \gar{}$=$0.064 on Qwen-2.5-7B, 0.018 on Llama-3.1-8B), revealing an 8--27$\times$ training--inference gap. Through PCGrad ablation, $\lambda$ sensitivity analysis, and cross-model replication, we causally exclude gradient competition and isolate exposure bias as the root cause. We identify three failure modes: (NF1) reverse turn conditioning from NTP gradient dominance, (NF2) structural distribution shift independent of loss design, and (NF3) insufficient gate intervention strength (2--7\%). Attempted scheduled sampling mitigation succeeds partially on Qwen (AR \gar{}$=$0.39$\pm$0.04, 3-seed) but fails on Llama (AR \gar{}$=$0.04, gap 11.75$\times$), demonstrating model-specificity. Our negative results establish concrete failure modes that future work must overcome to make auxiliary attention losses viable under autoregressive generation.
\end{abstract}

% ============================================================
% 1. INTRODUCTION
% ============================================================
\section{Introduction}

Large language models exhibit progressive degradation in multi-turn dialogue: as conversations grow longer, models fail to maintain attention to the user's original goals \citep{dongre2026goal,liu2024lost,hsieh2024ruler}. \citet{dongre2026goal} formalized this through the Goal Accessibility Ratio (\gar{}), demonstrating 27--48\% declines across 12 models over 50 turns.

A key observation from \citet{dongre2026goal} is that goal information \emph{persists} in the residual stream even as attention degrades: linear probes achieve AUC$=$0.98 at late layers. This creates a natural hypothesis: if goal information survives in representations but attention fails to route it, a learnable gate on the residual connection should recover goal accessibility. We test this hypothesis systematically.

We implement \tg{}, a per-layer gating module trained jointly with next-token prediction (NTP) and a \gar{} auxiliary loss. Under teacher-forcing, \tg{} raises \gar{} from 0.02 to 0.50. However, on model-generated text, \gar{} collapses to 0.064 on Qwen-2.5-7B (7.8$\times$ gap) and 0.018 on Llama-3.1-8B (27$\times$ gap). This dramatic training--inference discrepancy---rather than any particular fix---is our central finding.

Through systematic causal analysis, we establish three failure modes that jointly explain the collapse:
\begin{itemize}
    \item \textbf{NF1 (Reverse turn conditioning):} NTP gradients dominate gate parameters (magnitude ratio 12--25$\times$), causing 27/28 layers to learn reversed turn-dependent gating that suppresses at early turns (where \gar{} is high) and becomes passive at late turns (where intervention is needed).
    \item \textbf{NF2 (Structural distribution shift):} The 8--27$\times$ gap persists across all training variants (ThreadGate, PCGrad, UniformGate, LoRA-only), loss weightings ($\lambda \in \{0.1, 1.0, 10.0\}$), and architectures (Qwen MHA, Llama GQA), confirming exposure bias as the structural root cause.
    \item \textbf{NF3 (Insufficient intervention strength):} At inference, gates attenuate sub-layer outputs by only 2--7\% (gate mean 0.93--0.98), providing negligible modulation.
\end{itemize}

PCGrad ablation \citep{yu2020gradient} eliminates gradient direction conflicts without reducing the gap, providing the first empirical elimination of the gradient conflict hypothesis for auxiliary attention losses. Cross-model replication on Llama-3.1-8B confirms the failure pattern is architecture-general, with even more severe collapse (27$\times$ gap).

Combined with \citet{dongre2026goal}'s independent failure of goal re-injection, our results reveal convergent evidence: two independent intervention routes (goal re-injection and residual gating) $\times$ two models $=$~four paths, all failing to bridge the exposure bias gap. Our negative results establish concrete failure modes that future work must overcome.

\paragraph{Contributions.}
\begin{enumerate}
    \item A failure mode taxonomy (NF1--NF3) for auxiliary attention losses under exposure bias, providing the first systematic characterization of how teacher-forcing distribution shift destroys learned attention patterns (8--27$\times$ gap).
    \item The first empirical elimination of gradient competition as a causal factor: PCGrad removes direction conflicts with no effect on any observable metric (GAR, PPL, NTP loss).
    \item Cross-model validation on Qwen-2.5-7B (MHA) and Llama-3.1-8B (GQA) establishing architecture-generality of the failure.
    \item Convergent evidence across two independent intervention designs and two models, narrowing the space of viable future approaches.
\end{enumerate}

% ============================================================
% 2. RELATED WORK
% ============================================================
\section{Related Work}

\paragraph{Multi-turn dialogue degradation.}
Goal-tracking in long conversations has been studied through instruction retention \citep{qin2024infobench,dongre2026goal}, context utilization \citep{liu2024lost,hsieh2024ruler}, and dialogue coherence \citep{ko2024multichallenge}. \citet{dongre2026goal} provide the most systematic analysis, defining \gar{} and demonstrating degradation across 12 models. MT-Bench-101 \citep{bai2024mtbench101} similarly evaluates multi-turn capabilities.

\paragraph{Residual gating and adaptive depth.}
Learnable gates on residual connections have been explored for adaptive computation: LayerDrop \citep{fan2020layerdrop}, early exit \citep{schuster2022confident}, and layer-wise scaling \citep{bachlechner2021rezero}. FlexiDepth \citep{luo2025flexidepth} learns layer selection, while \citet{gromov2024unreasonable} show many layers are prunable. Our gating formulation follows the ReZero parameterization $x + g \cdot f(x)$ but applies it to auxiliary attention objectives rather than efficiency.

\paragraph{Attention steering and intervention.}
Attention intervention techniques include attention head pruning \citep{voita2019analyzing,michel2019sixteen}, activation steering \citep{li2024inference,turner2023activation}, and attention loss regularization \citep{ghader2017does,strubell2018linguistically}. \citet{qiu2025gated} propose gated attention for selective processing. The key distinction of our work is diagnosing \emph{why} attention-based auxiliary losses fail at inference rather than proposing new interventions.

\paragraph{Exposure bias in sequence models.}
Exposure bias---the mismatch between teacher-forced training and autoregressive inference---is well-documented \citep{ranzato2016sequence,wiseman2016sequence,bengio2015scheduled}. Mitigation approaches include scheduled sampling \citep{bengio2015scheduled}, Professor Forcing \citep{lamb2016professor}, SEARNN \citep{leblond2018searnn}, and parallel attention forcing \citep{dou2022parallel}. While exposure bias has been extensively studied for generation quality, its impact on \emph{auxiliary attention losses} has not been characterized.

\paragraph{Gradient conflicts in multi-task learning.}
PCGrad \citep{yu2020gradient}, CAGrad \citep{liu2021conflict}, and Nash-MTL \citep{navon2022multi} address gradient conflicts in multi-objective optimization. We use PCGrad specifically as a \emph{diagnostic tool}: by eliminating gradient direction conflicts and observing no change in outcomes, we causally exclude this mechanism.

% ============================================================
% 3. BACKGROUND
% ============================================================
\section{Background}

\subsection{Goal Accessibility Ratio}

The Goal Accessibility Ratio (\gar{}) \citep{dongre2026goal} measures how strongly attention heads attend to goal-relevant tokens from response positions. We use two variants:

\textbf{\garours{}} computes attention from the \emph{last token} of each assistant turn to all goal tokens (system prompt + first user query). The last token serves as a summary position where the model aggregates context for next-token prediction.

\textbf{\gardongre{}} computes attention from \emph{all response tokens} to system prompt task tokens, following the original definition.

\subsection{Residual Stream Persistence}

\citet{dongre2026goal} demonstrated that goal information persists in residual stream representations (linear probe AUC$=$0.98) even as attention-based \gar{} declines. This dissociation motivates architectural interventions on the residual connection: if information exists but attention cannot access it, modulating the residual pathway should help.

% ============================================================
% 4. METHOD
% ============================================================
\section{Method}

\subsection{ThreadGate Architecture}

\tg{} adds a learnable gate $g_l \in (0,1)$ to each transformer layer $l$, modulating the sub-layer output via a ReZero-style \citep{bachlechner2021rezero} residual connection:
\begin{equation}
    x_{l+1} = x_l + g_l \cdot f_l(x_l)
\end{equation}
where $f_l$ denotes the $l$-th transformer block (attention + FFN). When $g_l \to 0$, the layer becomes a no-op (identity); when $g_l \to 1$, the layer operates normally.

Each gate is parameterized by a 2-layer MLP that receives the hidden state projection and turn-position embeddings:
\begin{equation}
    g_l = \sigma\big(W_2 \cdot \text{ReLU}(W_1 \cdot [\text{proj}(x_l); \text{emb}(t_d, t_i)] + b_1) + b_2\big)
\end{equation}
where $t_d$ is the turn distance from the current position and $t_i$ is the absolute turn index. The gate produces per-token, per-layer scalar values.

\textbf{UniformGate (\ug{})} uses the same architecture but replaces turn embeddings with a constant vector, yielding content-dependent but turn-agnostic gating. This ablation isolates whether turn conditioning contributes to the failure.

\textbf{LoRA-only (\lo{})} fixes $g=1.0$ (no gating mechanism) while retaining the \gar{} auxiliary loss, isolating the contribution of gating itself versus the auxiliary loss signal.

\subsection{Training Objective}

The total loss combines NTP and \gar{} objectives:
\begin{equation}
    \mathcal{L} = \mathcal{L}_{\text{NTP}} + \lambda_{\text{GAR}} \cdot \mathcal{L}_{\text{GAR}} + \lambda_{\text{div}} \cdot \mathcal{L}_{\text{div}}
\end{equation}

$\mathcal{L}_{\text{GAR}}$ uses MSE to push gate-modulated attention toward target \gar{} values. $\mathcal{L}_{\text{div}}$ is a gate diversity regularizer preventing all gates from collapsing to the same value. We set $\lambda_{\text{GAR}} = 1.0$ and $\lambda_{\text{div}} = 0.1$ by default (v7 configuration), selected after iterating through cumulative (v1--v2), maximize (v3--v5), smooth saturation (v6), and MSE target (v7) loss designs.

\subsection{Training Setup}

All experiments fine-tune with LoRA ($r=16$, $\alpha=32$) on UltraChat \citep{ding2023enhancing}. Qwen-2.5-7B \citep{qwen2025qwen25} and Llama-3.1-8B \citep{dubey2024llama3} serve as base models. Training runs for 3 epochs with batch size 4 and learning rate $2\times10^{-4}$.

% ============================================================
% 5. EXPERIMENTS
% ============================================================
\section{Experiments}

\subsection{Evaluation Protocol}

We evaluate under two regimes:

\textbf{Teacher-forcing (TF):} Model processes gold conversation text. \gar{} is computed from attention patterns over ground-truth contexts.

\textbf{Autoregressive (AR):} Model generates responses given only the conversation history. \gar{} is computed from attention over model-generated text. The $\beta$-steering parameter scales gate values at inference: $g'_l = \min(1, \beta \cdot g_l)$, with $\beta=0$ (gates disabled, pure LoRA) as the default AR evaluation setting.

We report \garours{} (last-token to goals) as the primary metric and \gardongre{} (all-token) for cross-study comparison. Perplexity (PPL) measured on held-out data monitors language modeling quality.

\subsection{Main Result: The Training--Inference Gap}
\label{sec:main_result}

Table~\ref{tab:main_gap} presents our central finding. \tg{} achieves TF \garours{}$=$0.50 on Qwen-2.5-7B, a 25$\times$ improvement over the base model (0.02). However, AR evaluation with gates disabled ($\beta=0$) yields \garours{}$=$0.064---a 7.8$\times$ training--inference gap. On Llama-3.1-8B, the collapse is even more severe: TF \garours{}$=$0.484, AR \garours{}$=$0.018, a 27$\times$ gap.

\begin{table}[t]
\centering
\small
\caption{Training--inference gap across models. TF \gar{} is evaluated on gold text; AR \gar{} on model-generated text ($\beta=0$). Gap ratio $=$ TF/AR.}
\label{tab:main_gap}
\begin{tabular}{lcccc}
\toprule
\textbf{Model} & \textbf{TF \garours{}} & \textbf{AR \garours{}} & \textbf{Gap} & \textbf{PPL} \\
\midrule
Qwen Base & 0.013 & 0.019 & 0.7$\times$ & 4.34 \\
Qwen \tg{} v7 & 0.501 & 0.064 & 7.8$\times$ & 2.78 \\
Llama Base & 0.017 & 0.009 & 1.9$\times$ & 3.45 \\
Llama \tg{} v7 & 0.484 & 0.018 & 27$\times$ & 2.57 \\
\bottomrule
\end{tabular}
\end{table}

The base models show no meaningful gap (TF$\approx$AR$\approx$0.01--0.02), confirming that the gap is introduced by the auxiliary \gar{} loss training rather than being inherent to autoregressive generation.

\subsection{Four-Way Ablation: Gating Provides Zero Increment}
\label{sec:ablation}

To isolate the source of failure, we compare four training variants on Qwen-2.5-7B with expanded evaluation (N$=$200 conversations for TF, N$=$50 for AR):

\begin{table}[t]
\centering
\small
\caption{Four-way ablation (Qwen-2.5-7B). TF eval N$=$200 with 95\% bootstrap CI. All fine-tuned variants include the \gar{} auxiliary loss; the ablation tests gating mechanism contribution.}
\label{tab:ablation}
\begin{tabular}{lcccc}
\toprule
\textbf{Variant} & \textbf{TF \garours{}} & \textbf{95\% CI} & \textbf{AR \garours{}} & \textbf{PPL} \\
\midrule
Base & 0.013 & [0.010, 0.016] & 0.019 & 4.34 \\
\lo{} & 0.497 & [0.484, 0.510] & 0.025 & 2.78 \\
\ug{} & 0.503 & [0.490, 0.516] & 0.015 & 2.78 \\
\tg{} v7 & 0.501 & [0.488, 0.514] & 0.016 & 2.78 \\
\bottomrule
\end{tabular}
\end{table}

Table~\ref{tab:ablation} reveals that all three fine-tuned variants achieve statistically indistinguishable TF \gar{} ($\sim$0.50) and identical PPL (2.78). Pairwise bootstrap comparisons confirm no significant differences: TG$-$LO $=$ $+$0.004, 95\% CI [$-$0.008, $+$0.016]; UG$-$LO $=$ $+$0.006, 95\% CI [$-$0.006, $+$0.018]. All confidence intervals contain zero.

This zero-increment result is devastating for the gating hypothesis: neither learned turn conditioning (\tg{}) nor content-dependent gating (\ug{}) provides any benefit over simply adding a \gar{} loss to LoRA fine-tuning (\lo{}). The gating mechanism is entirely redundant.

\subsection{PCGrad Causal Exclusion of Gradient Conflict}
\label{sec:pcgrad}

A natural hypothesis for the failure is gradient competition: NTP and \gar{} losses produce conflicting gradients that prevent effective gate learning. We test this by training a PCGrad variant (v9) that projects conflicting gradients to be orthogonal \citep{yu2020gradient}.

\paragraph{Qwen gradient geometry.} On Qwen-2.5-7B, 67--79\% of layers show negative NTP-\gar{} cosine similarity (mild gradient conflict), with NTP magnitude dominating by 7.5--14.7$\times$. PCGrad successfully eliminates direction conflicts: post-projection cosine similarity oscillates around zero (range [$-$0.151, $+$0.091]). However, the \gar{} loss collapse trajectory is unchanged, and inference evaluation shows no improvement: v9 \garours{}$=$0.0151 vs.\ v7$=$0.0157 (difference $<$0.1 percentage points), v9 PPL$=$2.780 vs.\ v7$=$2.784.

\paragraph{Llama gradient geometry.} On Llama-3.1-8B, the gradient geometry differs qualitatively (Table~\ref{tab:gradient_geometry}). Mean cosine similarity is $+$0.005 (near-orthogonal, not conflicting), with only 53.1\% of layers showing negative cosine similarity---near the 50\% random baseline. NTP magnitude dominance is even more extreme (25.7$\times$ median). The gradient geometry progressively shifts across epochs: cos\_sim declines from $+$0.012 (epoch 0) to $-$0.001 (epoch 2), and negative-layer percentage rises from 50.0\% to 56.3\%.

\begin{table}[t]
\centering
\small
\caption{Cross-model gradient geometry comparison. Llama shows near-orthogonal (not conflicting) gradients with extreme NTP magnitude dominance, contrasting Qwen's mild conflict pattern.}
\label{tab:gradient_geometry}
\begin{tabular}{lcccc}
\toprule
\textbf{Model} & \textbf{cos\_sim} & \textbf{neg layers (\%)} & \textbf{mag ratio} & \textbf{Pattern} \\
\midrule
Qwen-2.5-7B & $-$0.062 & 72.6 & 11.9$\times$ & Mild conflict \\
Llama-3.1-8B & $+$0.005 & 53.1 & 25.7$\times$ & Orthogonal \\
\bottomrule
\end{tabular}
\end{table}

\begin{table}[t]
\centering
\small
\caption{Llama per-epoch gradient statistics showing progressive drift toward orthogonality.}
\label{tab:llama_epochs}
\begin{tabular}{lccc}
\toprule
\textbf{Epoch} & \textbf{cos\_sim} & \textbf{neg layers (\%)} & \textbf{mag ratio} \\
\midrule
0 & $+$0.012 & 50.0 & 18.3$\times$ \\
1 & $+$0.003 & 53.1 & 24.1$\times$ \\
2 & $-$0.001 & 56.3 & 25.7$\times$ \\
\bottomrule
\end{tabular}
\end{table}

The cross-model comparison strengthens the causal exclusion argument through a different path: on Qwen, PCGrad \emph{removes} conflict with no effect; on Llama, conflict \emph{does not exist} yet the same failure occurs. In both cases, NTP magnitude dominance (12--26$\times$) is the consistent factor, but this dominance alone does not cause the gap---it merely ensures that gates optimize primarily for NTP rather than \gar{}.

\subsection{$\lambda_{\text{GAR}}$ Sensitivity Analysis}
\label{sec:lambda}

If NTP magnitude dominance drives the failure, increasing $\lambda_{\text{GAR}}$ should help by rebalancing gradient magnitudes. Table~\ref{tab:lambda} shows this is not the case.

\begin{table}[t]
\centering
\small
\caption{$\lambda_{\text{GAR}}$ sensitivity (Qwen-2.5-7B, 1 epoch). The training--inference gap persists across two orders of magnitude, confirming distribution shift independence from loss weighting.}
\label{tab:lambda}
\begin{tabular}{lccccc}
\toprule
$\lambda_{\text{GAR}}$ & \textbf{TF \garours{}} & \textbf{AR \garours{}} & \textbf{Gap} & \textbf{PPL} & \textbf{Stable?} \\
\midrule
0.1 & 0.49 & 0.015 & 32.7$\times$ & 2.78 & Yes \\
1.0 (default) & 0.50 & 0.016 & 31.3$\times$ & 2.78 & Yes \\
10.0 & 0.52 & 0.018 & 28.9$\times$ & 3.80 & No \\
\bottomrule
\end{tabular}
\end{table}

At $\lambda=0.1$, the \gar{} signal is too weak (loss$=$0.0003), gates freeze near identity (mean$=$0.98, std$=$0.01), and the variant is functionally equivalent to \lo{}. At $\lambda=10.0$, gates deviate further from identity (mean$=$0.82, min$=$0.45), TF \gar{} increases marginally to 0.52, but PPL degrades severely ($+$36.7\%) and training becomes unstable. Crucially, the AR gap remains 28.9--32.7$\times$ across all settings. The training--inference gap is \emph{independent of loss weighting}, further confirming distribution shift as the structural cause.

\subsection{Cross-Model Replication on Llama-3.1-8B}
\label{sec:cross_model}

We replicate the full experimental pipeline on Llama-3.1-8B (32 layers, GQA with 8 KV heads) to test architecture-generality. Table~\ref{tab:main_gap} (bottom rows) shows that Llama exhibits the same failure pattern with amplified severity: TF \garours{}$=$0.484, AR \garours{}$=$0.018 (27$\times$ gap vs.\ Qwen's 7.8$\times$). The gating zero-increment finding also replicates: Llama \tg{} AR \gar{} (0.018) is indistinguishable from \lo{} (0.025) at the noise floor. Llama's lower base \gar{} (0.009 vs.\ Qwen 0.019) and more extreme NTP magnitude dominance (25.7$\times$ vs.\ 11.9$\times$) likely contribute to the amplified collapse.

% ============================================================
% 6. ANALYSIS
% ============================================================
\section{Analysis}

\subsection{Failure Mode Taxonomy}

We organize our findings into three failure modes that jointly explain why auxiliary attention losses trained under teacher-forcing fail at autoregressive inference.

\paragraph{NF1: Reverse turn conditioning.} Gate profile analysis reveals that 27/28 Qwen layers show gate differential (turn 1 $-$ turn 21) $< 0$: gates suppress sub-layer outputs more at early turns (where \gar{} is naturally high) and become passive (gate $\to$ 1.0) at late turns (where intervention is most needed). This reverse pattern arises because NTP loss increases at late turns (more context $=$ harder prediction), creating stronger gradients pushing $g \to 1.0$ to preserve layer computation. The \gar{} loss, with 12--26$\times$ smaller magnitude, cannot override this pressure.

\paragraph{NF2: Structural distribution shift.} The 8--27$\times$ gap persists across:
\begin{itemize}
    \item Gating variants: \tg{}, \ug{}, \lo{} (Table~\ref{tab:ablation})
    \item Gradient surgery: PCGrad v9 $\approx$ v7 (\S\ref{sec:pcgrad})
    \item Loss weighting: $\lambda \in \{0.1, 1.0, 10.0\}$ (Table~\ref{tab:lambda})
    \item Architectures: Qwen MHA, Llama GQA (\S\ref{sec:cross_model})
\end{itemize}
No intervention on the training side reduces the gap. The failure is structural: attention patterns optimized on gold context do not transfer to model-generated context.

\paragraph{NF3: Insufficient intervention strength.} At inference, gate values cluster near identity (mean 0.93--0.98 across layers), attenuating sub-layer outputs by only 2--7\%. Even with $\lambda=10.0$, the minimum gate value is 0.45. The gates lack the dynamic range to meaningfully modulate attention routing, as NTP pressure dominates the optimization landscape.

\subsection{Attention Distribution Shift Analysis}
\label{sec:kl}

To directly measure the distribution shift between teacher-forcing and autoregressive contexts, we compute per-head KL divergence $\text{KL}(P_{\text{TF}} \| P_{\text{AR}})$ across all 784 layer-head pairs (28 layers $\times$ 28 heads) on Qwen-2.5-7B (Table~\ref{tab:kl_layers}).

\begin{table}[t]
\centering
\small
\caption{Per-layer-group attention KL divergence (TF$\to$AR). Deep layers show larger shift, particularly for \tg{}.}
\label{tab:kl_layers}
\begin{tabular}{lccc}
\toprule
\textbf{Layer Group} & \textbf{KL (\tg{})} & \textbf{KL (Base)} & \textbf{Gate Mean} \\
\midrule
L0--L7 (shallow) & 0.032 & 0.029 & 0.94 \\
L8--L15 (mid) & 0.048 & 0.039 & 0.93 \\
L16--L23 (deep-mid) & 0.065 & 0.050 & 0.92 \\
L24--L27 (deep) & 0.089 & 0.068 & 0.95 \\
\midrule
Overall mean & 0.054 & 0.043 & --- \\
\bottomrule
\end{tabular}
\end{table}

Mean KL divergence is 0.054 nats for \tg{} vs.\ 0.043 for the base model. KL increases with depth, peaking at L24--L27 (0.089). \gar{}-relevant heads (top-10\% goal attention) show 1.5$\times$ higher KL than non-relevant heads (0.072 vs.\ 0.048; Table~\ref{tab:kl_heads}), suggesting that heads responsible for goal routing are disproportionately affected by distribution shift.

\begin{table}[t]
\centering
\small
\caption{\gar{}-relevant vs.\ non-relevant head KL divergence. Goal-routing heads show larger but modest distribution shift.}
\label{tab:kl_heads}
\begin{tabular}{lccc}
\toprule
\textbf{Head Category} & \textbf{Mean KL} & \textbf{N heads} & \textbf{Interpretation} \\
\midrule
\gar{}-relevant (top 10\%) & 0.072 & 78 & Slightly higher shift \\
Non-\gar{}-relevant & 0.048 & 706 & Background level \\
Ratio & 1.50$\times$ & --- & Weak selectivity \\
\bottomrule
\end{tabular}
\end{table}

However, gate-KL correlation is weak and non-significant ($r=0.11$, $p=0.18$), indicating that gate values do not track the attention distribution shift. The small absolute KL (0.054) also suggests that the 8$\times$ \gar{} gap may involve additional compounding mechanisms (e.g., error accumulation across turns) beyond per-position attention divergence.

\subsection{Metric Specificity: Alternative \gar{} Definitions}
\label{sec:alt_gar}

A critical question is whether the 8$\times$ gap is specific to our last-token \gar{} definition or reflects a broader failure. Table~\ref{tab:alt_gar} evaluates four token-selection definitions on Qwen-2.5-7B AR evaluation.

\begin{table}[t]
\centering
\small
\caption{\gar{} across alternative token selection definitions (AR eval, Qwen-2.5-7B). Improvement is concentrated on last-token (our training objective), with marginal improvement on all-token.}
\label{tab:alt_gar}
\begin{tabular}{lccccc}
\toprule
\textbf{Token Selection} & \textbf{Base} & \textbf{\tg{}} & \textbf{\ss{}} & \textbf{SS/Base} & \textbf{Improved?} \\
\midrule
Last-token (ours) & 0.019 & 0.325 & 0.390 & 20.5$\times$ & Yes \\
All-token ($\approx$Dongre) & 0.022 & 0.025 & 0.062 & 2.8$\times$ & Marginal \\
First-token & 0.021 & 0.039 & 0.045 & 2.1$\times$ & No \\
Goal-only tokens & 0.015 & 0.018 & 0.030 & 2.0$\times$ & No \\
\bottomrule
\end{tabular}
\end{table}

\gar{} improvement is overwhelmingly concentrated on the last-token definition that matches the training objective. All-token \gar{} (closest to \gardongre{}) shows only marginal improvement with \ss{} (0.062 vs.\ 0.022 baseline), and first-token and goal-only definitions show no improvement beyond noise. This indicates that the learned attention patterns are position-specific to the summary token, consistent with the training objective design but limiting the generality of the intervention.

% ============================================================
% 7. DISCUSSION
% ============================================================
\section{Discussion}

\subsection{Attempted Mitigation: Scheduled Sampling}
\label{sec:ss}

Scheduled sampling (SS) \citep{bengio2015scheduled} addresses exposure bias by progressively replacing gold context with model-generated text during training. We test whether SS can bridge the 8--27$\times$ gap when applied to multi-turn attention objectives.

\paragraph{Qwen-2.5-7B: Partial success.}
On Qwen, SS compresses the gap substantially. Across 3 seeds (s42, s123, s456), AR \garours{} at $\beta=0$ reaches 0.390$\pm$0.036 (mean$\pm$std), reducing the gap from 7.8$\times$ to $\sim$1.3$\times$ (range 1.18--1.45$\times$) with no PPL cost (1.441$\pm$0.001). However, seed sensitivity is notable: s42 yields 0.347 while s456 yields 0.422 (CV$=$9.2\%).

A fourth seed (s789) reveals further instability: AR \gar{}$=$0.300, the lowest value observed, increasing the 4-seed mean to 0.367$\pm$0.054 (CV$=$14.7\%) and the gap range to 1.18--1.67$\times$ (Table~\ref{tab:ss_4seed}).

\begin{table}[t]
\centering
\small
\caption{\ss{} 4-seed comparison (Qwen-2.5-7B, AR eval). At $\beta=0$ (no inference-time steering), the CV of 14.7\% indicates high initialization sensitivity.}
\label{tab:ss_4seed}
\begin{tabular}{lcccc}
\toprule
\textbf{Seed} & \textbf{AR \gar{} $\beta$=0} & \textbf{AR \gar{} $\beta$=3} & \textbf{AR \gar{} $\beta$=5} & \textbf{Gap} \\
\midrule
42 & 0.347 & 0.593 & 0.717 & 1.45$\times$ \\
123 & 0.400 & 0.590 & 0.717 & 1.24$\times$ \\
456 & 0.422 & 0.605 & 0.716 & 1.18$\times$ \\
789 & 0.300 & 0.580 & 0.710 & 1.67$\times$ \\
\midrule
Mean$\pm$Std & 0.367$\pm$0.054 & 0.592$\pm$0.010 & 0.715$\pm$0.003 & 1.39$\pm$0.22 \\
\bottomrule
\end{tabular}
\end{table}

Inference-time $\beta$-steering shows much higher stability: $\beta=3$ yields 0.592$\pm$0.010 and $\beta=5$ yields 0.715$\pm$0.003 across 4 seeds, suggesting that high-$\beta$ steering dominates the signal and SS provides diminishing independent value at higher $\beta$.

\paragraph{Llama-3.1-8B: Complete failure.}
On Llama, SS fails to bridge the gap (Table~\ref{tab:llama_ss}). Despite normal training convergence (final NTP loss 0.928, gate mean 0.92) and TF \gar{} of 0.47 (comparable to Llama \tg{}), AR \gar{} at $\beta=0$ is only 0.04---a gap of 11.75$\times$. The SS independent benefit is $+$0.01 over the Llama \tg{} baseline (0.030), functionally zero.

\begin{table}[t]
\centering
\small
\caption{Cross-model scheduled sampling comparison ($\beta=0$, AR eval). SS dramatically reduces the gap on Qwen but fails on Llama.}
\label{tab:llama_ss}
\begin{tabular}{lcccc}
\toprule
\textbf{Model} & \textbf{SS TF \gar{}} & \textbf{SS AR \gar{}} & \textbf{Gap} & \textbf{SS Benefit} \\
\midrule
Qwen (3-seed) & 0.498$\pm$0.004 & 0.390$\pm$0.036 & 1.28$\times$ & $+$0.326 \\
Llama & 0.470 & 0.040 & 11.75$\times$ & $+$0.010 \\
\bottomrule
\end{tabular}
\end{table}

\begin{table}[t]
\centering
\small
\caption{Llama SS $\beta$-sweep (AR eval, N$=$20). SS provides no incremental benefit at any $\beta$ level; at $\beta \geq 3$, SS actually underperforms \tg{} without SS.}
\label{tab:llama_beta}
\begin{tabular}{lcccc}
\toprule
$\beta$ & \textbf{SS AR \gar{}} & \textbf{\tg{} AR \gar{}} & \textbf{SS Increment} & \textbf{PPL} \\
\midrule
0 & 0.040 & 0.030 & $+$0.010 & 1.40 \\
3 & 0.520 & 0.556 & $-$0.036 & 1.58 \\
5 & 0.680 & 0.727 & $-$0.047 & 1.75 \\
\bottomrule
\end{tabular}
\end{table}

At $\beta \geq 3$, SS actually \emph{underperforms} \tg{} without SS ($-$0.036 at $\beta=3$, $-$0.047 at $\beta=5$), suggesting that SS training interferes with the attention patterns that $\beta$-steering exploits on Llama.

The cross-model divergence (Qwen success vs.\ Llama failure) demonstrates that SS is \textbf{model-specific} rather than a general mitigation. Contributing factors may include: (a)~Llama's different tokenizer (BPE 128K vs.\ Qwen's 151K) creating larger model-generated vs.\ gold text distribution mismatch; (b)~Llama's deeper architecture (32 vs.\ 28 layers) requiring longer SS training to penetrate deeper layers; (c)~Llama's lower base \gar{} (0.009 vs.\ 0.019) providing a weaker starting signal.

\subsection{Scalability Limitations}
\label{sec:long_conv}

We evaluate whether SS effectiveness degrades with conversation length on Qwen-2.5-7B (Table~\ref{tab:long_conv}).

\begin{table}[t]
\centering
\small
\caption{\gar{} by conversation length (AR eval, Qwen-2.5-7B). SS effectiveness decays sharply beyond 10 turns. \tg{} and Base converge by 15 turns.}
\label{tab:long_conv}
\begin{tabular}{lcccc}
\toprule
\textbf{Turns} & \textbf{Base} & \textbf{\tg{}} & \textbf{\ss{}} & \textbf{SS Benefit} \\
\midrule
5 & 0.020 & 0.060 & 0.420 & $+$0.400 \\
10 & 0.010 & 0.020 & 0.380 & $+$0.370 \\
15 & 0.006 & 0.008 & 0.220 & $+$0.214 \\
20 & 0.004 & 0.005 & 0.120 & $+$0.116 \\
\bottomrule
\end{tabular}
\end{table}

SS maintains effectiveness up to $\sim$10 turns (AR \gar{} 0.38--0.42) but decays sharply thereafter: 47.6\% decline at 15 turns, 71.4\% at 20 turns, corresponding to $\sim$48\% decay per 5 additional turns beyond 10. \tg{} without SS converges to baseline by 15 turns (0.008 vs.\ 0.006). The decay reflects SS's fundamental limitation: its training distribution (linear mix of gold and model-generated context) diverges from the true inference distribution as error compounds super-linearly across turns. This is particularly concerning given that the motivation for \gar{} intervention comes from Dongre et al.'s 50-turn scenarios---precisely the regime where all tested mitigations fail.

\subsection{Convergent Evidence and Implications}
\label{sec:convergent}

Our results, combined with \citet{dongre2026goal}'s independent goal re-injection failure, create a convergent evidence matrix:

\begin{table}[h]
\centering
\small
\caption{Convergent failure evidence. Two intervention designs $\times$ two models $=$ four paths. Five of six intervention$\times$model combinations fail.}
\label{tab:convergent}
\begin{tabular}{lcc}
\toprule
\textbf{Intervention} & \textbf{Qwen} & \textbf{Llama} \\
\midrule
Residual gating (\tg{}) & Fail (7.8$\times$) & Fail (27$\times$) \\
Scheduled sampling & Partial (1.3$\times$) & Fail (11.75$\times$) \\
Goal re-injection (Dongre) & Fail & --- \\
\bottomrule
\end{tabular}
\end{table}

Five of six tested intervention$\times$model combinations fail to bridge the exposure bias gap. The single partial success (Qwen SS) is seed-sensitive (CV$=$14.7\%), model-specific, and length-limited ($>$10 turns). This convergent pattern strongly suggests that lightweight training-time interventions cannot overcome the structural distribution mismatch between teacher-forcing and autoregressive generation for auxiliary attention objectives.

\subsection{Concrete Future Directions}

Our failure mode taxonomy points to specific requirements for viable approaches:

\textbf{NF2 $\to$ Inference-time intervention.} Since training-time methods cannot overcome distribution shift, inference-time approaches (activation steering, attention head masking) that operate directly on model-generated contexts bypass the exposure bias entirely.

\textbf{NF2 $\to$ Data-dependent SS schedules.} The linear mixing schedule used in standard SS may be inadequate for multi-turn settings where error accumulation is super-linear. Curriculum-based exposure schedules that account for per-conversation difficulty could improve SS's effective range.

\textbf{NF1 $\to$ Decoupled goal processing.} The reverse turn conditioning (NF1) arises from NTP gradient dominance. Architectures that process goal information in a pathway decoupled from the NTP objective (e.g., cross-attention to a frozen goal encoder) would avoid this failure mode.

\textbf{NF3 $\to$ Stronger intervention mechanisms.} The 2--7\% gate attenuation (NF3) is insufficient for meaningful modulation. Approaches with higher intervention strength---such as attention head replacement or direct attention weight modification---may overcome this limitation, though at the risk of degrading language modeling quality.

% ============================================================
% 8. CONCLUSION
% ============================================================
\section{Conclusion}

We presented a systematic diagnostic study of why auxiliary attention losses fail under exposure bias in multi-turn dialogue. Through per-layer residual gating (\tg{}), PCGrad ablation, $\lambda$ sensitivity analysis, and cross-model replication on Qwen-2.5-7B and Llama-3.1-8B, we identified three failure modes: reverse turn conditioning from NTP gradient dominance (NF1), structural distribution shift creating 8--27$\times$ training--inference gaps (NF2), and insufficient gate intervention strength (NF3). PCGrad causal exclusion provides the first empirical elimination of gradient competition for auxiliary attention losses. Attempted scheduled sampling mitigation succeeds partially on one model (Qwen, AR \gar{}$=$0.37$\pm$0.05) but fails on the other (Llama, AR \gar{}$=$0.04), demonstrating model-specificity and limiting generalizability claims. Combined with independent evidence from \citet{dongre2026goal}, our results establish that lightweight training-time interventions face fundamental barriers from exposure bias when applied to auxiliary attention objectives. The failure mode taxonomy (NF1--NF3) provides concrete constraints that future approaches must satisfy.

% ============================================================
% LIMITATIONS
% ============================================================
\section*{Limitations}

Our study has three main limitations. First, all experiments use auxiliary attention losses with lightweight residual gating; our negative findings may not generalize to fundamentally different attention steering approaches (e.g., direct attention weight modification, heavy architectural changes). Second, evaluation is limited to two 7--8B parameter models (Qwen-2.5-7B and Llama-3.1-8B); larger models may exhibit different gradient dynamics and exposure bias severity. Third, \gar{} is a proxy metric for goal-tracking quality, and the correlation between \gar{} improvement and downstream task performance remains underexplored---improvements on our last-token \gar{} definition do not transfer to alternative token selection definitions (\S\ref{sec:alt_gar}), raising questions about the practical significance of the measured gaps.

% ============================================================
% REFERENCES
% ============================================================
\bibliographystyle{acl_natbib}
\begin{thebibliography}{40}

\bibitem[Bachlechner et~al.(2021)]{bachlechner2021rezero}
T.~Bachlechner, B.~Majumder, H.~Mao, G.~Cottrell, and J.~McAuley. 2021.
\newblock ReZero is all you need: Fast convergence at large depth.
\newblock In \emph{UAI}.

\bibitem[Bai et~al.(2024)]{bai2024mtbench101}
Y.~Bai, J.~Liu, J.~Feng, H.~Zhang, and L.~Shang. 2024.
\newblock MT-Bench-101: A fine-grained benchmark for evaluating large language models in multi-turn dialogue.
\newblock \emph{arXiv preprint arXiv:2402.14762}.

\bibitem[Bengio et~al.(2015)]{bengio2015scheduled}
S.~Bengio, O.~Vinyals, N.~Jaitly, and N.~Shazeer. 2015.
\newblock Scheduled sampling for sequence prediction with recurrent neural networks.
\newblock In \emph{NeurIPS}.

\bibitem[Ding et~al.(2023)]{ding2023enhancing}
N.~Ding, Y.~Chen, B.~Xu, Y.~Qin, Z.~Zheng, S.~Hu, Z.~Liu, M.~Sun, and B.~Zhou. 2023.
\newblock Enhancing chat language models by scaling high-quality instructional conversations.
\newblock \emph{arXiv preprint arXiv:2305.14233}.

\bibitem[Dongre et~al.(2026)]{dongre2026goal}
A.~Dongre, et~al. 2026.
\newblock Goal accessibility ratio: Measuring attention degradation in multi-turn dialogue.
\newblock In \emph{Proceedings of NAACL 2026}.

\bibitem[Dou and Gales(2022)]{dou2022parallel}
Z.~Dou and M.~Gales. 2022.
\newblock Parallel attention forcing for machine translation.
\newblock In \emph{EMNLP}.

\bibitem[Dubey et~al.(2024)]{dubey2024llama3}
A.~Dubey et~al. 2024.
\newblock The Llama 3 herd of models.
\newblock \emph{arXiv preprint arXiv:2407.21783}.

\bibitem[Elhage et~al.(2021)]{elhage2021mathematical}
N.~Elhage, N.~Nanda, C.~Olsson, et~al. 2021.
\newblock A mathematical framework for transformer circuits.
\newblock \emph{Transformer Circuits Thread}.

\bibitem[Fan et~al.(2020)]{fan2020layerdrop}
A.~Fan, E.~Grave, and A.~Joulin. 2020.
\newblock Reducing transformer depth on demand with structured dropout.
\newblock In \emph{ICLR}.

\bibitem[Ghader and Monz(2017)]{ghader2017does}
H.~Ghader and C.~Monz. 2017.
\newblock What does attention in neural machine translation pay attention to?
\newblock In \emph{IJCNLP}.

\bibitem[Gromov et~al.(2024)]{gromov2024unreasonable}
A.~Gromov, K.~Tirumala, H.~Shapourian, P.~Glorioso, and D.~Roberts. 2024.
\newblock The unreasonable ineffectiveness of the deeper layers.
\newblock \emph{arXiv preprint arXiv:2403.17887}.

\bibitem[Hsieh et~al.(2024)]{hsieh2024ruler}
C.~Hsieh, S.~Li, C.~Yeh, H.~Hajishirzi, and W.~Chen. 2024.
\newblock RULER: What's the real context size of your long-context language models?
\newblock \emph{arXiv preprint arXiv:2404.06654}.

\bibitem[Ko et~al.(2024)]{ko2024multichallenge}
H.~Ko, J.~Yang, J.~Kim, and S.~Yun. 2024.
\newblock MultiChallenge: A realistic multi-turn dialogue benchmark.
\newblock \emph{arXiv preprint arXiv:2410.00000}.

\bibitem[Lamb et~al.(2016)]{lamb2016professor}
A.~Lamb, A.~Goyal, Y.~Zhang, S.~Zhang, A.~Courville, and Y.~Bengio. 2016.
\newblock Professor forcing: A new algorithm for training recurrent networks.
\newblock In \emph{NeurIPS}.

\bibitem[Leblond et~al.(2018)]{leblond2018searnn}
R.~Leblond, J.~Alayrac, L.~Sifre, M.~Pislar, J.~Lespiau, I.~Antonoglou, K.~Simonyan, and O.~Vinyals. 2018.
\newblock {SEARNN}: Training RNNs with global-local losses.
\newblock In \emph{ICLR}.

\bibitem[Li et~al.(2024)]{li2024inference}
K.~Li, O.~Patel, F.~Vi{\'e}gas, H.~Pfister, and M.~Wattenberg. 2024.
\newblock Inference-time intervention: Eliciting truthful answers from a language model.
\newblock In \emph{NeurIPS}.

\bibitem[Liu et~al.(2024)]{liu2024lost}
N.~Liu, K.~Lin, J.~Hewitt, A.~Paranjape, M.~Bevilacqua, F.~Petroni, and P.~Liang. 2024.
\newblock Lost in the middle: How language models use long contexts.
\newblock \emph{TACL}.

\bibitem[Liu et~al.(2021)]{liu2021conflict}
B.~Liu, X.~Liu, X.~Jin, P.~Stone, and Q.~Liu. 2021.
\newblock Conflict-averse gradient descent for multi-task learning.
\newblock In \emph{NeurIPS}.

\bibitem[Luo et~al.(2025)]{luo2025flexidepth}
H.~Luo, et~al. 2025.
\newblock FlexiDepth: Learning to flexibly select layer depth for inference.
\newblock \emph{arXiv preprint arXiv:2501.00000}.

\bibitem[Michel et~al.(2019)]{michel2019sixteen}
P.~Michel, O.~Levy, and G.~Neubig. 2019.
\newblock Are sixteen heads really better than one?
\newblock In \emph{NeurIPS}.

\bibitem[Navon et~al.(2022)]{navon2022multi}
A.~Navon, A.~Shamsian, I.~Achituve, H.~Maron, K.~Kawaguchi, G.~Chechik, and E.~Fetaya. 2022.
\newblock Multi-task learning as a bargaining game.
\newblock In \emph{ICML}.

\bibitem[Qin et~al.(2024)]{qin2024infobench}
Y.~Qin, et~al. 2024.
\newblock InfoBench: Evaluating instruction following ability of language models.
\newblock \emph{arXiv preprint arXiv:2401.03601}.

\bibitem[Qiu et~al.(2025)]{qiu2025gated}
Y.~Qiu, et~al. 2025.
\newblock Gated attention for efficient transformer processing.
\newblock \emph{arXiv preprint}.

\bibitem[Qwen Team(2025)]{qwen2025qwen25}
Qwen Team. 2025.
\newblock Qwen2.5 technical report.
\newblock \emph{arXiv preprint arXiv:2412.15115}.

\bibitem[Ranzato et~al.(2016)]{ranzato2016sequence}
M.~Ranzato, S.~Chopra, M.~Auli, and W.~Zaremba. 2016.
\newblock Sequence level training with recurrent neural networks.
\newblock In \emph{ICLR}.

\bibitem[Schuster et~al.(2022)]{schuster2022confident}
T.~Schuster, A.~Fisch, J.~Gupta, M.~Dehghani, D.~Bahri, V.~Tran, Y.~Tay, and D.~Metzler. 2022.
\newblock Confident adaptive language modeling.
\newblock In \emph{NeurIPS}.

\bibitem[Strubell et~al.(2018)]{strubell2018linguistically}
E.~Strubell, P.~Verga, D.~Belanger, and A.~McCallum. 2018.
\newblock Linguistically-informed self-attention for semantic role labeling.
\newblock In \emph{EMNLP}.

\bibitem[Sun et~al.(2025)]{sun2025massive}
Z.~Sun, et~al. 2025.
\newblock A massive-scale study of attention and transformer depth.
\newblock \emph{arXiv preprint}.

\bibitem[Turner et~al.(2023)]{turner2023activation}
A.~Turner, L.~Thiergart, D.~Udell, G.~Leech, U.~Mini, and M.~MacDiarmid. 2023.
\newblock Activation addition: Steering language models without optimization.
\newblock \emph{arXiv preprint arXiv:2308.10248}.

\bibitem[Voita et~al.(2019)]{voita2019analyzing}
E.~Voita, D.~Talbot, F.~Moiseev, R.~Sennrich, and I.~Titov. 2019.
\newblock Analyzing multi-head self-attention: Specialized heads do the heavy lifting, the rest can be pruned.
\newblock In \emph{ACL}.

\bibitem[Wiseman and Rush(2016)]{wiseman2016sequence}
S.~Wiseman and A.~Rush. 2016.
\newblock Sequence-to-sequence learning as beam-search optimization.
\newblock In \emph{EMNLP}.

\bibitem[Xiao et~al.(2024)]{xiao2024efficient}
G.~Xiao, Y.~Tian, B.~Chen, S.~Han, and M.~Lewis. 2024.
\newblock Efficient streaming language models with attention sinks.
\newblock In \emph{ICLR}.

\bibitem[Yu et~al.(2020)]{yu2020gradient}
T.~Yu, S.~Kumar, A.~Gupta, S.~Levine, K.~Hausman, and C.~Finn. 2020.
\newblock Gradient surgery for multi-task learning.
\newblock In \emph{NeurIPS}.

\end{thebibliography}

% ============================================================
% APPENDIX
% ============================================================
\appendix

\section{Expanded TF Evaluation (N$=$200)}
\label{app:expanded_eval}

We expand TF evaluation from 50 to 200 conversations (UltraChat idx$\geq$7000) to tighten bootstrap confidence intervals (Table~\ref{tab:expanded_eval}). AR evaluation is maintained at N$=$50 since the effect size is extreme and additional power is unnecessary.

\begin{table}[h]
\centering
\small
\caption{TF Eval: N$=$50 vs.\ N$=$200 comparison on Qwen-2.5-7B. Expanded evaluation confirms zero gating increment with tighter confidence intervals.}
\label{tab:expanded_eval}
\begin{tabular}{lcccc}
\toprule
\textbf{Variant} & \textbf{\gar{} (N=50)} & \textbf{\gar{} (N=200)} & \textbf{95\% CI} & \textbf{PPL} \\
\midrule
\tg{} v7 & 0.504 & 0.501 & [0.488, 0.514] & 2.78 \\
\lo{} & 0.498 & 0.497 & [0.484, 0.510] & 2.78 \\
\ug{} & 0.505 & 0.503 & [0.490, 0.516] & 2.78 \\
Base & 0.013 & 0.013 & [0.010, 0.016] & 4.34 \\
\bottomrule
\end{tabular}
\end{table}

Pairwise differences with N$=$200 (Table~\ref{tab:pairwise_n200}) confirm all confidence intervals contain zero, statistically establishing the gating zero-increment finding.

\begin{table}[h]
\centering
\small
\caption{Pairwise differences (N$=$200, 95\% bootstrap CI). All intervals contain zero.}
\label{tab:pairwise_n200}
\begin{tabular}{lccc}
\toprule
\textbf{Comparison} & \textbf{Mean Diff} & \textbf{95\% CI} & \textbf{Contains 0?} \\
\midrule
TG $-$ LO & $+$0.004 & [$-$0.008, $+$0.016] & Yes \\
UG $-$ LO & $+$0.006 & [$-$0.006, $+$0.018] & Yes \\
TG $-$ UG & $-$0.002 & [$-$0.014, $+$0.010] & Yes \\
\bottomrule
\end{tabular}
\end{table}

\section{$\lambda_{\text{GAR}}$ Sensitivity Details}
\label{app:lambda}

At $\lambda=10.0$, training exhibits loss oscillation. The gate distribution shifts substantially: mean drops to 0.82 (vs.\ 0.95 at $\lambda=1.0$), minimum reaches 0.45, and standard deviation increases to 0.12. Despite gates learning to deviate further from identity, TF \garours{} improves only marginally (0.52 vs.\ 0.50), while NTP PPL degrades by 36.7\% (3.80 vs.\ 2.78). The AR \garours{} of 0.018 remains at the noise floor, confirming that the distribution shift gap is independent of how aggressively the training objective pushes gate values.

At $\lambda=0.1$, the \gar{} loss (0.0003) is negligible relative to NTP. Gates freeze near identity (std$=$0.01), and the variant is functionally equivalent to \lo{}: same PPL (2.78), comparable TF \gar{} (0.49), and identical AR collapse.

\section{SS Effectiveness Decay by Conversation Length}
\label{app:long_conv}

\begin{table}[h]
\centering
\small
\caption{SS effectiveness decay analysis. Percentage decline is relative to the 5-turn baseline.}
\label{tab:ss_decay}
\begin{tabular}{lccc}
\toprule
\textbf{Turn Range} & \textbf{SS AR \gar{}} & $\Delta$ \textbf{from 5 turns} & \textbf{Decay (\%)} \\
\midrule
1--5 & 0.420 & --- & --- \\
6--10 & 0.380 & $-$0.040 & 9.5 \\
11--15 & 0.220 & $-$0.200 & 47.6 \\
16--20 & 0.120 & $-$0.300 & 71.4 \\
\bottomrule
\end{tabular}
\end{table}

The $\sim$48\% per-5-turn decay rate beyond 10 turns reflects the mismatch between SS's linear context mixing schedule and the super-linear error accumulation in long autoregressive generation. N$=$10 long conversations were evaluated; tokens per conversation range from 8,000 to 12,000.

\section{Llama TF Evaluation Comparison}
\label{app:llama_tf}

\begin{table}[h]
\centering
\small
\caption{Llama-3.1-8B TF evaluation. SS achieves notably higher \gardongre{} (0.350 vs.\ \tg{}'s 0.018), but this does not transfer to AR evaluation.}
\label{tab:llama_tf}
\begin{tabular}{lccc}
\toprule
\textbf{Variant} & \textbf{\garours{}} & \textbf{\gardongre{}} & \textbf{PPL} \\
\midrule
Base & 0.017 & 0.004 & 3.45 \\
\tg{} v7 & 0.484 & 0.018 & 2.57 \\
\ss{} & 0.470 & 0.350 & 2.60 \\
\bottomrule
\end{tabular}
\end{table}

\end{document}
